\par\Needspace{6\baselineskip}\medskip
\begin{center}\hypertarget{native-178-article131}{}\pdfbookmark[2]{131.}{article-131}131.\end{center}
In art. 114 ostensum est, si ratio inter quantitates illic atque in art. 128
per $n,n',n''$ denotatas cognita fuerit, corporis coelestis distantias a terra per
formulas persimplices determinari posse. Quodsi itaque pro $x,y$ assumerentur
quotientes $\dfrac{n}{n'},\ \dfrac{n''}{n'}$, pro his quantitatibus in eo casu, ubi motus heliocentricus
inter observationes haud ita magnus est, statim valores approximati $\dfrac{\theta}{\theta'},\ \dfrac{\theta''}{\theta'}$
se offerent (accipiendo characteres $\theta,\theta',\theta''$ in eadem significatione ut in art. 128):
hinc itaque solutio obvia problematis nostri demanare videtur, si ex $x$ et $y$
distantiae duae a terra eliciantur, ac dein ad instar alicuius ex quinque methodis
artt. 124—128 procedatur. Revera, acceptis quoque characteribus $\eta,\eta''$
in significatione art. 128, designatoque analogice per $\eta'$ quotiente orto ex divisione
sectoris inter duos radios vectores contenti per aream trianguli inter
eosdem, erit
\[
\frac{n}{n'}=\frac{\theta}{\theta'}\cdot\frac{\eta'}{\eta}\mathpunct{\raisebox{.45ex}{,}}\qquad
\frac{n''}{n'}=\frac{\theta''}{\theta'}\cdot\frac{\eta'}{\eta''}\mathpunct{\raisebox{.45ex}{,}}
\]
patetque facile, si $n,n',n''$ tamquam quantitates parvae primi ordinis spectentur,
esse generaliter loquendo $\eta-1,\eta'-1,\eta''-1$ quantitates secundi ordinis,
adeoque valores ipsarum $x,y$ approximatos $\dfrac{\theta}{\theta'},\ \dfrac{\theta''}{\theta'}$ a veris differre tantummodo
quantitatibus secundi ordinis. Nihilominus re propius considerata methodus
haecce omnino inepta invenitur, cuius phaenomeni rationem paucis explicabimus.
Levi scilicet negotio perspicitur, quantitatem $(0.1.2)$, per quam
distantiae in formulis 9, 10, 11 art. 114 multiplicatae sunt, ad minimum tertii
ordinis fieri, contra e.\ g. in aequ. 9 quantitates $(\mathrm{O}.1.2)$, $(\mathrm{I}.1.2)$, $(\mathrm{II}.1.2)$ primi
ordinis; hinc autem facile sequitur, errorem secundi ordinis in valoribus quantitatum
$\dfrac{n}{n'},\ \dfrac{n''}{n'}$ commissum producere in valoribus distantiarum errorem ordinis
$0$. Quamobrem, secundum vulgarem loquendi usum, distantiae tunc quoque
errore finito affectae prodirent, quando temporum intervalla infinite parva
sunt, adeoque neque has distantias neque reliquas quantitates inde derivandas
ne pro approximatis quidem habere liceret, methodusque conditioni secundae
art. praec. adversaretur.
